\chapter{The Buddhist (Tibetan) Model}

Tibetan Buddhism is a fully worked-out account of mind, world, and path that is organized around the refusal of a maker. It inherits the Indian Buddhist model and adds the Vajrayana machinery of deity, mantra, mandala, the subtle body, and the after-death states. Everything is sorted into interlocking numbered systems. It is the most structurally elaborate cosmology in this collection.

\section*{No maker, no first moment}

The defining feature is stated at the outset. There is \textbf{no creator God} and there is \textbf{no first cause}. Every other tradition here retains a source from which things come. Tibetan Buddhism does not. Time is \textbf{beginningless}. The cosmos was not made and will not end. It cycles without a first moment. The absence of a maker is asserted as deliberately as another tradition would assert a maker.

\section*{What plays the role of the ultimate}

Even without a creator, the tradition has candidates for what is most fundamental. The schools disagree on which is final, and that disagreement is a live theme of Tibetan thought.

\begin{itemize}
 \item \textbf{Emptiness} (\textbf{shunyata}). The ultimate nature of everything is the lack of inherent existence. Not a thing and not a void, but the absence of any independent self-existence. The Gelug school and the Prasangika Madhyamaka hold this as final.
 \item \textbf{The nature of mind} (\textbf{rigpa}, the \textbf{Ground}). In Dzogchen the ultimate is the Ground, pure awareness, primordially pure and spontaneously present, the clear light. Not a creator, but the basis of both confusion and awakening.
 \item \textbf{Buddha-nature} (\textbf{tathagatagarbha}). The awakened potential present in all beings. The Jonang shentong view treats it as a truly existent luminous ground, against the Gelug reading of it as mere emptiness.
\end{itemize}

The ultimate is emptiness, or the nature of mind, or buddha-nature, depending on the school. None of these is a being. None of them creates. All of them are already present, to be recognized rather than made.

\section*{The four schools and the key texts}

The model is shared across the tradition. The emphasis differs by school.

\begin{itemize}
 \item \textbf{Nyingma}, the old school, home of Dzogchen, the Ground and pure awareness.
 \item \textbf{Kagyu}, home of Mahamudra and the Six Yogas of Naropa.
 \item \textbf{Sakya}, home of the path-and-fruit teaching, Lamdre.
 \item \textbf{Gelug}, home of the graded path Lamrim, Prasangika emptiness, and the strict self-emptiness view.
\end{itemize}

The \textbf{Jonang} school carries the other-emptiness view that the Gelug reject. The key texts run from the Abhidharmakosha of Vasubandhu, through the Prajnaparamita and Madhyamaka of Nagarjuna and Chandrakirti, the Yogachara of Asanga and Vasubandhu, the Bardo Thodol known in English as the Tibetan Book of the Dead, the Kalachakra Tantra, the Lamrim of Tsongkhapa, and the Dzogchen tantras.
